\section{Implementation}

This chapter !!!.

\subsection{Implementation Overview}
The simulator consists of 3 main parts: the trace decoder, the tag controller algorithm, and the logical tag cache. Figure \ref{fig:wbx} shows how these parts fit into the information flow of the whole project.

% what about the visualisation tool??? similar sentence for this?

[FIGURE WBX!!!]

\subsection{Simulator Framework}
- collect command line args
- opens trace files
- has decoder decode trace files
- opens output file and log file
- initialises controller (and thus tag cache state)
- begin simulation: iterate through llc accesses and have the tag controller handle them
- report number of resulting accesses and number of llc miss entries processed

\subsection{Trace Decoder}
- mention assumption scheme due to incomplete info
- lets main read in llc misses (uses both files, uses furnish->getInitial...)
- lets controller initialise its cache's state

\subsection{Tag Controller}
- has a logical cache
- abstract class's constructor initialises cache
- implementations must set the state of the cache (different schemes have different tables)
- handleMemoryAccess must is abstract
- some cache proxy methods

\subsubsection{Basic Scheme}
- simple cache setup
- simple translation

\subsubsection{\textit{Efficient Tagged Memory}'s Hierarchical Scheme}
- implementation of the algorithm (give code snippet and refer back to figure)

\subsubsection{The Morello Project Scheme}
- same cache setup as basic
- same translation as basic
- emulate caches seperately!

\subsection{Tag Cache}
- log methods (which are now INTERNAL; not used by other classes) log to the trace
- outward-facing do methods
- outward-facing peek and set; peek used by morello to see state without logging; set used in cache initialisation and also in morello to set without logging

\subsection{Logging}

\subsubsection{Logging Tag Table State}
- simulator is told the log points
- tells controller to dump the cache, the cache dumps to an output file
- format of the dump
- also return the points where champsim needs to log (i.e. after how many accesses, as it's not 1:1 with llc misses)

\subsubsection{Adding Logging to ChampSim for Tag Cache State}
- cache.h: made BLOCK public
- main.cc: new logpoints commandlinearg, open logfile in setup, pass it to champsim::main
- champsim.cc: after each round of update, check if the next logpoint has been reached, if so log to outputfile

\subsection{Visualisation Tool}
- parse command line args
- run controller sim, collecting log points for champsim
- print out number of accesses and the logpoints
- run champsim
- open the output logs
- draw things (in a senible order)
- save image, which must be opened by a tool such as vipsdisp

\subsection{Testing}

\subsection{Repository Overview}
- the dirtree thing
